\documentclass[12pt,a4paper]{article}

\usepackage[utf8]{inputenc}
\usepackage[T1]{fontenc}
\usepackage[margin=2.5cm]{geometry}
\usepackage{parskip}
\usepackage{enumitem}
\usepackage{xcolor}
\usepackage{mdframed}

\newenvironment{stagenote}
  {\begin{mdframed}[linecolor=blue!50, backgroundcolor=blue!5, linewidth=1pt]%
   \small\color{blue!70!black}\textbf{Stage direction:} }
  {\end{mdframed}}

\newenvironment{slidenote}
  {\begin{mdframed}[linecolor=green!50!black, backgroundcolor=green!5, linewidth=1pt]%
   \small\color{green!40!black}\textbf{Slide:} }
  {\end{mdframed}}

\title{Presentation Script\\[0.3em]\large The Gaia Mission --- Mapping the Milky Way}
\author{Michael Nedzelsky}
\date{22 March 2026, 11:00\,AM}

\begin{document}
\maketitle

\begin{stagenote}
Before the presentation: use a compass app on your phone to find
which direction is South in the room. Also check the app ``Sky Map'' (Android)
or ``Sky Guide'' (iPhone) to verify where the Sun, the Galactic Centre,
and Andromeda are at 11:00\,AM. Approximate directions at 11:00 GMT
from Cambridge on 22 March:
\begin{itemize}[nosep]
    \item Sun: south-southeast (azimuth $\sim$162\textdegree), about 37\textdegree{} above the horizon
    \item Galactic Centre: west-southwest (azimuth $\sim$245\textdegree), about 17\textdegree{} BELOW the horizon
    \item Andromeda Galaxy: east-southeast (azimuth $\sim$115\textdegree), about 71\textdegree{} up --- nearly overhead
\end{itemize}
\end{stagenote}

\section*{1. Introduction (1 minute)}

\begin{slidenote}
Title slide.
\end{slidenote}

Hello everyone. My name is Michael, and I work at the Institute of Astronomy here in Cambridge. Today I want to tell you about my project --- a space mission called Gaia. But first, let me remind you just how big the universe is, because I think it is easy to forget.

\section*{2. Where Are We Right Now? (3 minutes)}

\begin{slidenote}
``Where Are We Right Now?'' slide.
\end{slidenote}

Right now, we are sitting in this room in Cambridge. But where is Cambridge in the universe? Let me show you.

\begin{stagenote}
Point in the direction of the Sun. If there is a window on that side,
even better --- ``the sunlight is coming from there.''
\end{stagenote}

First --- the Sun. Right now, the Sun is in \emph{that} direction, roughly south. It is about 150 million kilometres away. Light from the Sun takes about 8 minutes to reach us. So the sunlight you see right now left the Sun 8 minutes ago.

Now, something more surprising. The centre of our galaxy --- the Milky Way --- where is it?

\begin{stagenote}
Point toward WSW and slightly downward, through the floor.
\end{stagenote}

It is in \emph{that} direction. Below the floor. The centre of our galaxy is below the horizon right now, about 26,000 light-years away. One light-year is about 10 trillion kilometres. So the centre of our galaxy is incredibly far.

And the nearest large galaxy to us? The Andromeda Galaxy?

\begin{stagenote}
Point nearly straight up, slightly toward ESE.
\end{stagenote}

It is almost directly above us. About 2.5 million light-years away. You can see a photograph of it on this slide --- a beautiful spiral galaxy, the nearest large one to us.

So right now, while we sit here, we are surrounded by the universe in every direction --- the galactic centre below us, a giant galaxy above us.

\begin{slidenote}
``How Big Is Andromeda in the Sky?'' slide.
\end{slidenote}

And here is something surprising. This image shows Andromeda next to the Moon at the same scale. You can see that Andromeda is actually several times wider than the full Moon on the sky. But its light is so faint that you cannot see it with your eyes, even though it is enormous.

\section*{3. How Far? (3 minutes)}

\begin{slidenote}
``How Far? --- The Solar System'' slide.
\end{slidenote}

These numbers --- millions, billions, trillions of kilometres --- are impossible to imagine. So let me use a comparison.

Imagine the Sun is the size of a basketball. A normal basketball, about 24 centimetres across. How big is the Earth at this scale?

\begin{stagenote}
Pause for guesses from the audience.
\end{stagenote}

The Earth would be a peppercorn. About 2 millimetres. And how far from the basketball? Twenty-six metres --- about the length of this street outside.

Jupiter, the biggest planet, would be the size of a walnut --- 134 metres from the basketball. Neptune would be a peanut, 774 metres away. And Pluto, which is no longer considered a planet, would be a grain of sand about one kilometre from the basketball.

Voyager~1, the most distant human-made object, launched in 1977, would be about 4 kilometres away at this scale. It has been travelling for almost 50 years and it is still inside our tiny neighbourhood.

And the nearest star? The nearest star would be \emph{another basketball}\ldots{} about 6,900 kilometres from here. Roughly the distance from Cambridge to Mumbai. That is how empty space is.

\begin{slidenote}
``How Far? --- The Galaxy'' slide.
\end{slidenote}

Now let me change the scale completely. If we shrink the entire Solar System --- Sun, all planets, everything out to Neptune --- to the size of a small coin, about 2 centimetres\ldots{} then our galaxy, the Milky Way, would be about the size of Cambridge. About 20 kilometres across.

And the Andromeda Galaxy, the one that is above us right now? At this scale it would be near Edinburgh --- about 500 kilometres away.

And inside our Cambridge-sized galaxy, there are about 100 billion of these little coins. Each one is a star with its own planets. Most of the space between them is empty.

\section*{4. The Milky Way (1 minute)}

\begin{slidenote}
``The Milky Way --- Our Home Galaxy'' slide.
\end{slidenote}

Our galaxy is a spiral galaxy. We cannot photograph it from outside, of course --- we are inside it. But on this slide you can see a real photograph of a galaxy called UGC\,12158, which astronomers believe looks very similar to the Milky Way. You can see the beautiful spiral arms. Our Sun would be in one of these arms, about two thirds of the way out from the centre.

The Sun orbits around the centre of the galaxy, but very slowly. One orbit takes about 230 million years. The last time the Sun was in this position, dinosaurs were just appearing on Earth.

\section*{5. Star Catalogues Through History (2 minutes)}

\begin{slidenote}
``Star Catalogues Through History'' slide.
\end{slidenote}

People have been making catalogues of stars for thousands of years. A star catalogue is basically a list --- for each star, you write down where it is on the sky, and how bright it is. It sounds simple, but it is extremely important for astronomy and navigation.

The first serious catalogue was made by Hipparchus, a Greek astronomer, around 129 BC. He recorded about 850 stars. For about 1,500 years, catalogues like his were the best we had. Astronomers and sailors used them for centuries.

Then came telescopes. Tycho Brahe measured about 1,000 stars with great accuracy. Flamsteed at Greenwich recorded about 3,000. But still --- just a few thousand stars.

The real leap came in the 19th century, with photography. The Bonner Durchmusterung, a German catalogue from 1859, recorded over 324,000 stars. Suddenly, hundreds of thousands of stars were catalogued. The Henry Draper catalogue in the early 20th century added spectral types for 225,000 stars.

Then in 1989, the European Space Agency launched a satellite called Hipparcos --- named after Hipparchus. From space, away from the atmosphere, Hipparcos measured 118,000 stars with very high precision. A related catalogue, Tycho-2, reached 2.5 million stars.

And then came Gaia.

\section*{6. The Gaia Mission (3 minutes)}

\begin{slidenote}
``What Is Gaia?'' slide.
\end{slidenote}

Gaia is a space telescope that was launched by the European Space Agency in December 2013. Its goal was to create the most detailed three-dimensional map of our galaxy. The satellite operated for 11 years --- more than double the planned five years, and took its last observation in January 2025.

What does ``three-dimensional'' mean here? Earlier catalogues told you \emph{where} a star is on the sky --- like a flat photograph. Gaia also measured \emph{how far away} each star is. And \emph{how fast} it is moving, and \emph{in which direction}. So Gaia gives us a complete picture --- positions, distances, and motions of stars.

One more fun fact: Gaia had the largest digital camera ever sent to space --- 106 sensors with nearly one billion pixels.

\begin{slidenote}
``Lagrange Point L2'' slide.
\end{slidenote}

Now, where was Gaia? It was not in orbit around the Earth like the Space Station. It was much further away, at a special place called Lagrange Point 2, or L2.

L2 is a point about 1.5 million kilometres from Earth, on the opposite side from the Sun. At this point, the gravity of the Sun and the gravity of the Earth work together to keep a satellite in a stable position. Gaia stayed behind the Earth, so the Sun, Earth, and Moon were all on one side --- behind a sunshield. This means Gaia had a clear, dark, stable view of the sky.

By the way, the James Webb Space Telescope is also at L2, right now.

Going back to our basketball model --- if the Sun is a basketball and the Earth is a peppercorn 26 metres away, then L2 is about 26 centimetres behind the peppercorn. Very close to Earth, in cosmic terms.

\begin{slidenote}
``How Gaia Scans the Sky'' slide.
\end{slidenote}

Gaia had two telescopes pointing in two different directions, separated by about 106 degrees. The satellite slowly spun --- one full turn every 6 hours. As it spun, the two telescopes scanned a strip of the sky. And the direction of the spin axis itself slowly changed over about 63 days. This combination of spinning and changing direction means that, over 11 years, Gaia covered the entire sky many times. Each star was observed about 70 times or more.

\section*{7. Results (1.5 minutes)}

\begin{slidenote}
``Gaia's Results'' slide.
\end{slidenote}

So what has Gaia found? The latest big data release, called Data Release 3, was published in June 2022. It contains positions of about 1.8 billion stars. For about 1.5 billion of them, we also have distances and motions.

To give you a sense of scale: Hipparchus catalogued 850 stars. Gaia catalogued nearly 2 billion. That is more than two million times more stars.

And it is not just stars. Gaia also measured asteroids, distant galaxies, and quasars. Over 12,000 scientific papers have been published using Gaia data.

The satellite stopped observing last year, but the work is not finished. Data Release~4, using all the data from the full mission, is expected later this year. And the final Data Release~5 is expected around 2030. Processing this much data takes years.

This is what I work on at the Institute of Astronomy --- helping to process and analyse this enormous amount of data.

\section*{8. Legacy (1 minute)}

\begin{slidenote}
``The Legacy of Star Catalogues'' slide.
\end{slidenote}

I want to end with something that gives me a lot of satisfaction.

Remember Hipparchus, who made his catalogue in 129 BC? That catalogue was still being used 1,500 years later. Tycho Brahe's observations from around 1600 were used by Kepler to discover the laws of planetary motion --- laws we still use today. The Hipparcos catalogue from 1997 is still used by astronomers, nearly 30 years later.

Star catalogues are special because they do not go out of date. The stars move, but very slowly. A good star catalogue remains useful for decades, sometimes for centuries.

Gaia's catalogue of 2 billion stars will be used by astronomers long after all of us in this room are gone. There is something very satisfying about working on something that will last.

\begin{stagenote}
Pause. Smile.
\end{stagenote}

Thank you. Do you have any questions?

\begin{slidenote}
``Thank you / Questions'' slide.
\end{slidenote}

\section*{Notes for preparation}

\begin{itemize}
    \item \textbf{Total speaking time:} approximately 14--15 minutes at a comfortable pace.
    \item \textbf{Compass app:} Install ``Sky Map'' (Android) or ``Sky Guide'' (iOS) to verify directions before the presentation. Also good to show the audience interactively.
    \item \textbf{Props (optional):} A basketball and a peppercorn would make the scale comparison very memorable.
    \item \textbf{Vocabulary to practise:}
    \begin{itemize}
        \item catalogue (British spelling) --- KAT-uh-log
        \item Hipparchus --- hih-PAR-kus
        \item parallax --- PAR-uh-laks
        \item Lagrange --- luh-GRAHNZH (French origin)
        \item quasar --- KWAY-zar
    \end{itemize}
    \item \textbf{Images included in slides:} Andromeda (Wikimedia), Andromeda+Moon (APOD), UGC\,12158 (ESA/Hubble), Gaia satellite (ESA), Lagrange points (NASA), Gaia sky map (ESA/Gaia/DPAC)
\end{itemize}

\end{document}
